% XVI-P p14/p14r science note: the P.5 ENTRY contract (quantum object, one prescription + D11 matching note, counterterm
% domain, executable amplitude rules). No loop is evaluated. Numbers are read from certificates/P14_*.json at build time
% (manuscript/body/P14-entry-summary.tex, P14-dossier-table.tex).
\section{Scope of revision p14 (repaired in p14r)}
\scopebox{p14r = a narrow repair of p14; P.3/P.4 code byte-identical; no loop evaluated.}{%
\begin{itemize}
\item The prescription is made operational: executable DRED index rules, tests against premature $d=4$, and the
withdrawn claim that no evanescent remainder arises (Section~\ref{sec:p14r-dred}).
\item The D11 M3 labels are corrected, and M3 is made conditional on M2, with an explicit normalization condition.
\item The counterterm quotient is worded as candidate directions before the identity constraints.
\item The regeneration coverage is stated, and the independent Hessian-route regeneration W6 is added for D01, D02, D03
and D12.
\item The isolated-copy gate tests contain the frozen inputs.
\end{itemize}
p14 is archived in \texttt{archival/p14/}.}
\scopebox{p14 = P.5 entry: the contract is fixed; no loop is evaluated.}{%
p14 fixes what P.5 will compute and how: the quantum object $\Gamma_\Psi$ in transformed variables
(Section~\ref{sec:p14-object}); one prescription together with a D11 matching note against XVI-O's frozen D11
(Section~\ref{sec:p14-presc}); the one-loop counterterm domain at $\kappa^1,\kappa^2$ (Section~\ref{sec:p14-ct}); and
executable amplitude rules with symbolic integrands for D01, D02, D03, D11 (reference only), D12, D13, D14a--c
(Section~\ref{sec:p14-amp}). No integral is evaluated: the Batch A/B/C integrals belong to p15 and later.
Opening P.5 entry earns none of the loop or restoration flags. \texttt{quantum\_inputs\_evaluated},
\texttt{projected\_defect\_computed}, \texttt{projected\_identity\_verified} and \texttt{projected\_restoration\_verified}
are not earned; \texttt{bulk\_symmetry\_restored}=false. D12 is Missing, never zero. O.1 is NOT\_CLOSED. The insertion
stays frozen, XV-B and XIV-J are not opened, Route B is not discharged, and no $c\,\delta$ is added to $u^R_\mu$. XVI-P
stays \emph{unregistered}: the package is evidence in the sense of XVI-O \texttt{conv:package}, not a series certificate,
and XVI-A--XVI-O are not edited.}

\section{P.3/P.4 frozen as P.5 inputs}\label{sec:p14-freeze}
An independent rebuild of p13 passed. A reviewer unpacked \texttt{XVI-P-source-only.zip} (sha256 \texttt{2f85a566\ldots})
in a fresh environment and ran the source-only entrypoint. All reference values matched (112/112, no mismatch), 77
values were recomputed by separate code, and a hand-written spot check of the conformal boundary identity agreed.
The report is archived as \texttt{archival/p13/reproducibility/independent\_rebuild\_report\_p13.json} (sha256 pinned).
\textbf{Caveat:} the rebuild ran on the same box (shared kernel, TeX Live and poppler) in a fresh virtual environment
from the source-only zip; it was not on a second machine.

The flag is \emph{bound to the p13 source tree} \texttt{2b1baf33\ldots} and to the p13 zip; p14 changes cannot inherit
it, and p14 has its own (pending) independent flag. \texttt{P3\_P4\_frozen\_as\_P5\_inputs} is recomputed in every build
by \texttt{freeze\_p13.py} as the conjunction of three conditions:
F1, the archived independent report has its pinned hash, \texttt{pass}=true, no mismatch, and is bound to the p13 tree and zip;
F2, the 54 frozen files (classical P.3/P.4 code, the inherited XVI-O code and notes, pins and inherited data;
\texttt{archival/p13/frozen\_p13\_inputs.json}) are byte-identical;
F3, the regenerated classical values (export hashes, the seven classical certificates and two ledgers, residual
records, pins) equal \texttt{archival/p13/reference\_values\_p13.json}.
The classical code is therefore unchanged in p14. Its certificates keep their \texttt{P13\_*} names, and the p14 gate
recomputes the p13 acceptance from them (Section~\ref{sec:frozen-classical}). \textbf{P.5 entry is stated against this
freeze:} if any of F1--F3 fails, P.5 entry is not earned and publication is rejected.

\section{Milestones (p14 status)}\label{sec:p14-milestones}
\begin{center}\small
\begin{tabular}{@{}p{2.6cm}p{9.2cm}p{3.4cm}@{}}
\toprule
milestone & definition (unchanged since p13, \texttt{conventions.MILESTONES}) & status (p14)\\
\midrule
P.3 & fixed output projection + exhaustive coefficient dependencies + graph topology inventory & complete (frozen)\\
P.4 & required classical action, transformations, source terms, gauge map and master verification.
\textbf{P.4 means classical readiness, NOT coupled one-loop restoration.} & earned (frozen)\\
P.5 entry & quantum prescription, amplitude rules and counterterm domain fixed & earned iff the five contract flags
hold and P.3/P.4 are frozen (see summary)\\
P.5 calculation & coefficients evaluated and renormalized & open (p15+)\\
P.5 identity & projected defect computed and, if possible, cancelled & open\\
\bottomrule
\end{tabular}
\end{center}
P.5 entry := \texttt{quantum\_action\_defined} $\wedge$ \texttt{quantum\_prescription\_defined} $\wedge$
\texttt{D11\_matching\_defined} $\wedge$ \texttt{counterterm\_domain\_defined} $\wedge$ \texttt{amplitude\_rules\_verified},
each recomputed from its certificate, together with \texttt{P3\_P4\_frozen\_as\_P5\_inputs}. The values of this build:
\begin{center}\small
\begin{tabular}{@{}p{5.2cm}p{10.3cm}@{}}
\toprule contract item & value (from the certificates of this build)\\ \midrule
integrated fields & h\_(mn) (10 independent components), c\^{}mu, cbar\_mu, phi, b\^{}mu (retained) \\
external antifields & K'\_(mn): h*' (10 comps, sigma\_h = +1), c*'\_mu: c*', cbar*'\^{}mu: cbar*', phi*': phi*', b*': absent (s b = 0) \\
b & RETAINED and integrated; det M\_bb = -16; eliminating b gives 4 terms quadratic in cbar*'; Gamma\_1 independent of bs, cbs \\
antifield shift & eps = -1 (derived) \\
prescription & P14\_DRED\_MSbar\_FeynmanI0: MS-bar: subtract the pole in the combination 1/eps - gamma\_E + log 4 pi; scale mu \\
i0 & Feynman: every propagator denominator q\^{}2 + i0 (eta = +,-,-,-); Wick rotation l0 = i l4 implied \\
DRED rules & verified = true; T1 finite a/8 + b/4 (premature d=4 would give b/4); T2 finite -a/2 ((d-4) x pole); mutations: premature\_d4 rejected, trace\_ghat\_eq\_4 rejected, contract\_after\_laurent rejected \\
D11 matching & defined = true, verified = false; M1: 100 component pairs, 0 mismatches; M3 (conditional on M2): kernel sd \textless{}= 8, contact order \textless{}= 4; contact space order\_0 2, order\_2 4, order\_4 5; under N4 5 (transverse 2); M2 defined\_not\_verified \\
counterterm domain & kappa\^{}1: 0; kappa\^{}2: 11 (h.h 5, c.cb 2, c.K 3, phi.phi 1); sigma\_0-equivalent: 3 (h.h 2, c.cb 0, c.K 0, phi.phi 1); 8 candidate directions modulo the computed exact equivalences, before the identity constraints (D cap ker sigma\_0 not computed) \\
amplitude controls & W1\_wick\_berezin\_shadow: pass, W2\_multiplicity: pass, W3\_closed\_form\_vertices: pass, W5\_regeneration: pass, W6\_hessian\_route: pass, W4\_ward\_bose: pass; W5 regenerates D14a-c, D11, D13; W6 (Hessian route) regenerates D01, D02, D03, D12 == exports: true \\
mutations (must be rejected) & wrong\_tadpole\_factor: rejected, wrong\_ghost\_loop\_sign: rejected, offdiag\_multiplicity: rejected, reversed\_source\_ghost\_order: rejected, omit\_ghost\_transport\_term: rejected \\
P.5 entry flags & quantum\_action\_defined = true, quantum\_prescription\_defined = true, D11\_matching\_defined = true, counterterm\_domain\_defined = true, amplitude\_rules\_verified = true \\
P3\_P4\_frozen\_as\_P5\_inputs & true (the p13 independent rebuild ran on the SAME box (shared kernel, TeX Live, poppler) in a fresh venv from the source-only zip; it was NOT run on a second machine) \\
P.5 entry & earned (contract fixed; stated against the P.3/P.4 freeze) \\
\bottomrule
\end{tabular}
\end{center}

\begin{center}\small
\begin{tabular}{@{}lllrrrl@{}}
\toprule graph & prefactor & denominators & comps & monomials & deg & status \\ \midrule
D01 & \texttt{kappa**2/2} & l\^{}2+i0, (l-p)\^{}2+i0 & 100 & 3888 & 4 & integrand\_available \\
D02 & \texttt{-kappa**2/2} & l\^{}2+i0 & 94 & 220 & 2 & integrand\_available \\
D03 & \texttt{-kappa**2} & (l-p)\^{}2+i0, l\^{}2+i0 & 100 & 2728 & 4 & integrand\_available \\
D11 & \texttt{kappa**2/2} & l\^{}2+i0, (l-p)\^{}2+i0 & 100 & 1722 & 4 & reference \\
D12 & \texttt{-kappa**2/2} & l\^{}2+i0 & 94 & 160 & 2 & Missing \\
D13 & \texttt{I*kappa**2} & l\^{}2+i0, (p-l)\^{}2+i0 & 40 & 396 & 3 & integrand\_available \\
D14a & \texttt{-kappa/2} & l\^{}2+i0 & 10 & 22 & 2 & integrand\_available \\
D14b & \texttt{kappa} & l\^{}2+i0 & 4 & 16 & 2 & integrand\_available \\
D14c & \texttt{-kappa/2} & l\^{}2+i0 & 10 & 22 & 2 & integrand\_available \\
\bottomrule
\end{tabular}
\end{center}


\section{The quantum object}\label{sec:p14-object}
\noindent Builder/output: \texttt{quantum\_object.py}, \texttt{P14\_quantum\_object.json}.\par\smallskip
\textbf{Action.} The action is computed directly in transformed variables:
\[S_\Psi(\Phi,\Phi^{*\prime})=\sum_k\kappa^{k-2}\!\int\ell_k+S_\phi+X\big(\Phi,\Phi^{*\prime}+\varepsilon\,\partial\Psi/\partial\Phi\big),
\qquad\varepsilon=-1\ (\text{derived}),\]
with $\Psi=-i\int\bar c_\mu(2D^\mu[h]-b^\mu)$. By order, $X$ contains $K'c,\ \bar c c,\ b^2,\ bh,\ b\bar c^{*\prime}$ at
$\kappa^0$; $K'ch,\ \bar cch,\ c^{*\prime}cc,\ \phi^{*\prime}c\phi$ at $\kappa^1$; and nothing at $\kappa^2$. $S_\phi$
contributes $\phi\phi$, $h\phi\phi$, $hh\phi\phi$. There are no switch or insertion fields: the baseline is $S_{\rm inv}$,
not $S_D$.

\textbf{Fields and sources.} The integrated fields are $h_{(mn)}$ (10 independent components), $c^\mu$, $\bar c_\mu$,
$\phi$ and $b^\mu$. The external antifields are $K'_{(mn)}$ ($\sigma_h=+1$), $c^{*\prime}_\mu$, $\bar c^{*\prime\mu}$ and
$\phi^{*\prime}$; $b^{*\prime}$ is absent because $sb=0$.

\textbf{$b$ is retained and integrated.} Four reasons:
\begin{itemize}
\item $b$ occurs only in $\kappa^0$ bilinears, so it has no interaction vertex;
\item $M_{bb}$ is constant and invertible ($\det=-16$);
\item the full $(h,b)$ inverse is verified, with Schur complement $q^2P$;
\item eliminating $b$ (computed) would make $S_\Psi$ \emph{quadratic} in $\bar c^{*\prime}$, whereas retaining it keeps
$S_\Psi$ linear in every external antifield.
\end{itemize}
$\bar c^{*\prime}$ couples only through $\bar c^{*\prime}b$, which has one quantum leg and is tree-level and
one-particle reducible. So the 1PI $\Gamma_1$ is independent of $\bar c^{*\prime}$ and $b^{*\prime}$, while $K'$,
$c^{*\prime}$ and $\phi^{*\prime}$ enter through $K'ch$, $c^{*\prime}cc$ and $\phi^{*\prime}c\phi$.

\textbf{1PI convention.} The weight is $e^{iS_\Psi/\hbar}$ and $\Gamma_\Psi=S_\Psi+\hbar\Gamma_1+O(\hbar^2)$, with the
Legendre transform taken in the fields only (antifields are spectators). $\Gamma_1=\tfrac i2{\rm STr}\log S^{(2)}_\Psi
=-i\log Z|_{\rm 1PI}$, so that $i\Gamma_1$ is the sum of 1PI one-loop graphs (Gate~0; the Berezin/Wick shadow below).

\textbf{Vacuum and one-point.} The vacuum is normalised by $\Gamma_1[0]=0$. The one-point function (D14a--c) is retained:
no tadpole cancellation and no background shift are imposed, and no value is assigned.

\textbf{Fourier, propagators and i0.} We use $\partial\to ik$ with incoming momenta. The propagators are
$\langle hh\rangle=i(P^{-1})_{ab}/(q^2+i0)$ (de Donder; independent components),
$\langle c^\mu\bar c_\rho\rangle=i(M_0^{-1})_{\mu\rho}=-\delta_{\mu\rho}/(q^2+i0)$ and $\langle\phi\phi\rangle=i/(q^2+i0)$,
all of them inverses of the verified Hessians; the $i0$ is Feynman throughout.

\textbf{Bookkeeping.} Each $n$-leg vertex is $O(\kappa^{n-2})$. At one loop, $\Gamma_1$ at $\kappa^k$ has exactly $k$
external legs ($\sum_v(n_v-2)=E$), and this is checked against the prefactor of every exported integrand.

\textbf{Pullback, not theorem.} $\Gamma_\Psi(\Phi,\Phi^{*\prime})=\Gamma(\Phi,\Phi^{*\prime}-\partial\Psi/\partial\Phi)$ is used
only as the classical pullback (verified at tree level). No quantum gauge-change theorem is asserted.

\section{Prescription and the D11 matching note}\label{sec:p14-presc}
\noindent Builder/output: \texttt{prescription\_d11.py}, \texttt{P14\_prescription\_D11\_matching.json}.\par\smallskip
\textbf{One prescription.} Dimensional reduction (DRED), made operational in p14r
(Section~\ref{sec:p14r-dred}). Subtraction is $\overline{\rm MS}$ at scale $\mu$, and the $i0$ is Feynman. The rules for
scaleless integrals belong to the definition but are \emph{not applied}: no tadpole (D02, D12, D14a--c) receives a value.
D12 stays Missing. XVI-O fixed D12 and the one-point function by
\texttt{dimensional\_regularization\_m0}; P's prescription belongs to the same family, but P does not import that value.

\subsection{Operational DRED rules (p14r)}\label{sec:p14r-dred}
\noindent Builder/output: \texttt{dred.py}, \texttt{P14\_dred\_rules.json}; \texttt{quantum\_prescription\_defined}
requires \texttt{dred\_rules\_verified}, and the release gate recomputes it.\par\smallskip
\begin{description}
\item[R1, spaces.] The quasi-4-dimensional space Q4S splits as $\mathrm{Q}d\mathrm{S}\oplus\mathrm{Q}(4{-}d)\mathrm{S}$:
$g^{(4)}=\hat g+\tilde g$, with $\operatorname{tr}g^{(4)}=4$, $\operatorname{tr}\hat g=d$ and $\operatorname{tr}\tilde g=4-d$.
The metrics are projectors: $\hat g\hat g=\hat g$, $\tilde g\tilde g=\tilde g$, $\hat g\tilde g=0$, $g^{(4)}X=X$.
\item[R2, momenta.] The loop momentum $l$ \emph{and} the external momenta $p$ lie in Q$d$S ($p\cdot\tilde g=0$; a stated
choice). Field, vertex and polarisation indices are Q4S indices.
\item[R3, lifting.] The exported numerators are polynomials in the lower components $l_r,p_r$ ($r=0..3$) at fixed external
components. They are lifted as $l_r:=l\cdot e_{(r)}=l\cdot\hat e_{(r)}$ and $p_r=p\cdot e_{(r)}$, with $e_{(r)}$ the
$\eta$-orthonormal basis of Q4S. The vertex tensors are lifted with $\eta\to g^{(4)}$, so nothing is truncated.
\item[R4, reduction.] $\int l^{\mu_1}\cdots l^{\mu_n}F(l^2,l\cdot p)=\sum_kC_kB_k$, with $B_k$ the symmetric tensors built
from $\hat g$ and $p$, and $C_k$ given by the exact Gram (Passarino--Veltman) projection in symbolic $d$. Contracting with
the lifted basis gives $\hat g_{rs}=\eta_{rs}-\tilde G_{rs}$, where the evanescent $\tilde G_{rs}$ satisfy
$\eta^{rs}\tilde G_{rs}=4-d$ and $\tilde G_{rs}p^s=0$. The checked closed forms are:
the rank-2 tadpole $\hat g\,J/d$; the rank-4 tadpole $J/(d(d+2))$;
$B_{00}=(J_{10}p^2-J_{02})/((d-1)p^2)$; and $B_{11}=(dJ_{02}-p^2J_{10})/((d-1)p^4)$.
\item[R5, order.] \emph{Every} index contraction, including the external 4D projection, is done \emph{before} the
Laurent expansion. $d$ stays symbolic through the $\overline{\rm MS}$ subtraction and the finite-part extraction, so
$(d-4)\times$pole terms are retained.
\item[R6, evanescent.] Evanescent remainders \emph{do} arise in general. The lift of the exported integrands has
$\tilde G$ parts in every graph, and contraction produces $(4-d)\times$pole finite terms. They are kept. The p14 sentence
claiming that none arise is withdrawn.
\end{description}
\textbf{Tests (generic placeholders $J=a/\epsilon+b+c\epsilon$, $d=4-2\epsilon$; no integral of this theory is evaluated).}
\begin{itemize}
\item T1, the $\eta_{00}$ part of $\int l_0l_0F$, has finite part $a/8+b/4$. Substituting $d=4$ before the expansion
would give $b/4$; the guarded pipeline raises \texttt{PrematureD4Substitution}.
\item T2, $\eta^{rs}\times$ the $\tilde G$ part, equals $-(4-d)J/d$ with finite part $-a/2$, a $(d-4)\times$pole term.
\end{itemize}
The mutations $\operatorname{tr}\hat g=4$ (rules check fails; T1 becomes $b/4$ and T2 becomes 0) and
expand-before-contract (\texttt{UncontractedObjects}) are both rejected. Background literature (nothing is imported from
it): W.~Siegel, Phys.\ Lett.\ B84 (1979) 193 and B94 (1980) 37; D.~St\"ockinger, JHEP 0503 (2005) 076, hep-ph/0503129;
A.~Signer and D.~St\"ockinger, Nucl.\ Phys.\ B808 (2009) 88, arXiv:0807.4424.

\textbf{D11 matching (defined, not verified).}
\begin{description}
\item[M0, baseline.] $S_{\rm inv}$, not $S_D$. $a_1=\tfrac12\tau\eta-h$ is XVI-O's vertex (\texttt{lem:d11-prefactor}:
$\Pi=(i\kappa^2/4)T$).
\item[M1, exact integrand identity (computed).] For all 100 pairs $(A,B)$ of independent components,
\[-i\,\mathcal P\,\frac{N^{P}_{AB}(l,p)}{(l^2+i0)((l-p)^2+i0)}=\frac{i\kappa^2}{4}\,R_A^{\mu\nu}R_B^{\rho\sigma}
V^{ab}_{\mu\nu}V^{cd}_{\rho\sigma}\big[\tilde C_{ac}(l)\tilde C_{bd}(l-p)+\tilde C_{ad}(l)\tilde C_{bc}(l-p)\big],\]
where:
\begin{itemize}
\item $\tilde C_{ab}(k)=k_ak_b\,i/(k^2+i0)$ is the Fourier transform of $C=-\partial\partial D_F$;
\item $V$ is XVI-O's own \texttt{V\_tensor}, imported from \texttt{code/matter\_loop/d11\_tensor.py};
\item $R_A$ is the $\eta$-raising map $h^{\mu\nu}=\eta^{\mu a}\eta^{\nu b}h_{ab}$, with orbit multiplicity 2 for
off-diagonal $A$;
\item $\mathcal P=\kappa^2/2$ is P's prefactor, which carries the symmetry factor $\tfrac12$.
\end{itemize}
The identity matches $\Gamma^{(2)}_1=\tfrac12\int h^{\mu\nu}\Pi_{\mu\nu\rho\sigma}h^{\rho\sigma}$ (XVI-O) with P's ordered
extraction, and it is bound to the exported D11 integrand by its sha256.
\item[M2, nonlocal agreement on $U=\mathbb R^4\setminus\{0\}$ (defined only).]
$\langle\mathcal F^{-1}[P^{\rm DR}_{AB}],f\rangle=\langle R_AR_B\Pi^R_{\rm tr},f\rangle$ for $f\in\mathcal D(U)$, with
$\Pi^R_{\rm tr}=(i\kappa^2/4)\mathcal K u^R_\mu$, $\alpha=A^2/60$ and $\beta=A^2/18$ (\texttt{lem:BV\_rules\_on\_U},
all 256 components). This needs the evaluated amplitude, which is p15+ work.
\item[M3, local difference (conditional on M2).] $L_{\rm match}=\mathcal F^{-1}[P^{\rm DR}]-R R\,\Pi^R_{\rm tr}$.
The degrees are as follows. The scalar extension $u^R_\mu$ has scaling degree 4. The kernel $\mathcal Ku^R_\mu$
(with $\mathcal K$ of order 4) has scaling degree at most 8. A contact term has $\mathrm{sd}(\partial^\alpha\delta)=4+|\alpha|$.
\emph{Kernel scaling degree at most eight; degree of divergence at most four; hence the local ambiguity has derivative
order at most four.}

\emph{If} M2 holds, then $L_{\rm match}$ is supported at $\{0\}$ and lies in the bounded contact space of Lorentz-covariant
even $hh$ structures of derivative order $\le4$. This space is enumerated: order 0 ($\mu^4$) has dimension 2 (transverse 0);
order 2 ($\mu^2p^2$) has dimension 4 (transverse 1, the linearised Einstein--Hilbert structure); order 4 has dimension 5
(transverse 2). Odd orders cannot occur with four indices built from $p$ and $\eta$.

Exactly degree four is an \emph{explicit normalization condition} N4: no $\mu^4$ and no $\mu^2p^2$ contact. N4 is
\emph{not} implied by ``no other scale'', since $\mu$ is a scale. Under N4 the space is the 5-dimensional order-4 span,
equal to the $hh$ sector of the counterterm domain. XVI-O's contact ansatz is a candidate list
(\texttt{NOT\_proved\_spanning}) with three named order-4 entries, one of them the collective \texttt{L4\_mixed}; its
lower-order entries are excluded by N4. P's list is the explicit one.
\item[M4, finite adjustments.] Under M2 and N4, these are booked only as P-side counterterms in the $hh$ sector (without N4
the order-0/2 contacts would have to be added there). $u^R_\mu$ is not modified
(no $c\,\delta$), so the scalar-extension prescription \texttt{analytic\_FP\_u\_mu\_R\_o70} is preserved.
\item[M5, residue.] $r=-i\pi^2$ is an XVI-O-side statement, valid only \emph{after} the constructive holomorphy of
$F(\zeta)=(Q-i0)^\zeta$ (\texttt{lem:analytic\_FP\_u\_mu\_R\_o70}). P's $1/\epsilon$ pole is P-side, and no pole--residue
identification is asserted.
\item[M6/M7.] The XVI-O o70 files are byte-pinned. XVI-O's kernel Ward identity $\partial(\Pi_w+L)=0$ is
\emph{not} functional ST4b, and the projected identity (XVI-N \texttt{lem:field-only}) is distinct from bulk restoration.
\end{description}
Hence \texttt{D11\_matching\_defined}=true and \texttt{D11\_matching\_verified}=false.

\section{Counterterm domain}\label{sec:p14-ct}
\noindent Builder/output: \texttt{counterterm\_domain.py}, \texttt{P14\_counterterm\_domain.json}.\par\smallskip
This is an \emph{order-bounded EFT constraint list, not an admission} (XIV-J and XV-B are not opened). The one-loop
$\kappa^n$ part has $n$ external legs and local dimension $4+n$, with $[\kappa]=-1$; $[h]=[\phi]=[c]=[\bar c]=1$; and
$[K']=[c^{*\prime}]=[\phi^{*\prime}]=[\bar c^{*\prime}]=[b]=2$. The functionals have ghost number 0, are even under
$\phi\to-\phi$ (with $\phi^{*\prime}\to-\phi^{*\prime}$), and are Lorentz covariant (parity-odd structures vanish for
two-leg terms). The permitted species are the legs of loop-capable vertices; $b$, $\bar c^{*\prime}$ and $b^{*\prime}$
are excluded. The structures are \emph{generated} mechanically from $\eta$ and the leg momentum and turned into
densities, and their ranks are exact. Total derivatives vanish identically in the two-point representation.
\begin{itemize}
\item $\kappa^1$: empty (a one-leg term with derivatives is a total derivative, and there is no dimension-5 field).
\item $\kappa^2$: 11 structures.
\begin{itemize}
\item $hh$, 5: $h\Box^2h$, $\tau\Box^2\tau$, $(\partial\partial h)\Box\tau$, $(\partial h)\Box(\partial h)$,
$(\partial\partial h)^2$;
\item $\bar cc$, 2: $\bar c_\mu\Box^2c^\mu$, $\bar c_\mu\partial^\mu\Box\partial_\nu c^\nu$;
\item $K'c$, 3: $K'^{mn}\partial_m\partial_n\partial_rc^r$, $K'^m{}_m\Box\partial_rc^r$, $K'^{mn}\Box\partial_mc_n$;
\item $\phi\phi$, 1: $\phi\Box^2\phi$.
\end{itemize}
\end{itemize}
\textbf{Source dependence.} $K'$ enters only via $K'c$. $c^{*\prime}$ and $\phi^{*\prime}$ appear first in three-leg terms
at $\kappa^3$, outside the order bound.

\textbf{Equivalences.} Apply $\sigma_0=(S^{(0)}_\Psi,\cdot)$ (frozen antibracket) to every generated ghost-number $-1$,
dimension-5 bilinear $\Xi$. Antifield-linear $\Xi$ are field redefinitions (canonical transformations); field-only $\Xi$
are $s$-exact. The domain $\cap\operatorname{Im}\sigma_0$ is given in the summary; the $S^{(0)}_\Psi$
density used for this is checked against the verified Hessians. \emph{Eight candidate directions modulo the computed exact
equivalences, before imposing the relevant identity constraints.} This is not a cohomological classification, which
needs $D\cap\ker\sigma_0$ (the $\sigma_0$-closedness constraints) and the quotient
$(D\cap\ker\sigma_0)/(D\cap\operatorname{Im}\sigma_0)$. $D\cap\ker\sigma_0$ has not been computed, so the count 8 is an
upper bound only. The enumeration itself is unchanged.

\textbf{Normalisation.} $\overline{\rm MS}$ (pure poles at $\mu$). There is no $\kappa^2$ two-derivative counterterm, so
neither $\kappa$ nor the $h$ normalisation is renormalised at this order. This is the \emph{imposed} normalization
condition N4 of M3, not a consequence of masslessness. There is no $\kappa^1$ counterterm.
Finite adjustments stay inside the domain.

\section{Executable amplitude rules}\label{sec:p14-amp}
\noindent Builder/output: \texttt{feynman.py}, \texttt{shadow.py}, \texttt{amplitude\_rules.py},
\texttt{P14\_amplitude\_rules.json}, \texttt{exports/P14\_vertex\_tensors.json}, \texttt{exports/P14\_integrands/*.json}.\par\smallskip
\textbf{Rules.}
\begin{itemize}
\item \emph{Extraction:} ordered, $V[L_1..L_n]=\partial^L_{L_n}\cdots\partial^L_{L_1}S$, with leg 1 leftmost; the Feynman
vertex is $i\kappa^nV$.
\item \emph{Graviton legs:} independent components; the multiplicity (1 diagonal, 2 off-diagonal) enters only the tensor form.
\item \emph{Ghost:} $\bar c$ is on the left in $\bar cch$ and in $M_0$; a closed ghost loop carries $-1$.
\item \emph{Source:} $K'$ sits left of $c$ in $K'ch$ and in the $\Gamma_1$ coefficient $K'_{a'}G_{a'\nu}c^\nu$.
\item \emph{Symmetry factors:} $S_f=\tfrac12$ for D01, D02, D11, D12, D14a, D14c and $S_f=1$ for D03, D13, D14b.
\item \emph{Ghost-loop sign:} $-1$ for D03 and D14b.
\end{itemize}
The vertex tensors include every derivative term of the ghost-graviton vertex ($\bar c\,c\,h=2ia\,D^\rho[\mathcal L_ch]$,
including the transport term $c^r\partial_rh$) and the $K'$ vertex.

\textbf{Controls.}
\begin{description}
\item[W1.] Berezin/Wick shadow theory: the symmetry factors and Grassmann signs are fixed independently.
\item[W2.] A tensor-route recomputation of the multiplicities.
\item[W3.] Closed forms of the vertices.
\item[W4.] The primary evidence is the generic symbolic \emph{off-shell} Ward identities at $\kappa^1$ and $\kappa^2$.
A supplementary check is the ``Ward contraction with two free on-shell graviton legs and an off-shell gauge leg''. W4
also checks Bose symmetry.
\item[W5.] Independent sympy regeneration of D14a, D14b, D14c, D11 and D13 \emph{only}, plus a Bose-symmetry check of D01.
\item[W6 (p14r).] Independent regeneration of D01, D02, D03 and D12 from the one-loop expansion of the functional
(super)trace log of the field-dependent Hessian of $S_\Psi$ (\texttt{hessian\_route.py}), not from the Feynman-rule
assembler:
\[i\Gamma^{AB}=-\tfrac12\operatorname{Tr}(GU_2^{AB})+\tfrac12\operatorname{Tr}(GU_1^AGU_1^B)\ \text{(bosons)},\qquad
i\Gamma^{AB}=-\operatorname{Tr}(GU_1^AGU_1^B)\ \text{(ghosts)}.\]
The ingredients are an own monomial differentiation of $\ell_3,\ell_4$, an own inverse of the full $14\times14$ $(h,b)$
Hessian, and the ghost and scalar Hessians from the $S_\Psi$ densities. Structural checks: $X^{(1)}$ has no field-only
term and $X^{(2)}=0$, so there is no $\bar cchh$ vertex. The route equals the exported integrands exactly, component by
component (D01 100, D02 94, D03 100, D12 94; 0 mismatches), and it also rejects the tadpole-factor, ghost-sign,
multiplicity and transport-term mutations.
\end{description}
\textbf{Mutations that must be rejected (all rejected):} wrong tadpole factor, wrong ghost-loop sign, off-diagonal
multiplicity, reversed source/ghost order, and omitted ghost transport term. D13 has priority; its integrand is exported
with prefactor $i\kappa^2$ and order $(K',c)$. Integrands are symbolic (not integrated); the per-graph table is in the
summary above.

\textbf{Isolated-copy gate tests (p14r).} The test copies contain every frozen input (the top-level directories of
\texttt{frozen\_p13\_inputs.json}, including \texttt{tests/}); the p14 copies lacked \texttt{tests/}, so every copy was
unfrozen. The expected outcomes are:
\begin{itemize}
\item the unmutated copy has freeze true, and publication is accepted;
\item each P.5-only mutation (DRED mutation not rejected, counterterm flag, amplitude mutation) keeps the freeze
\emph{true}, turns the named contract flag false, and is rejected;
\item the hard-stop mutations keep the freeze true and are rejected;
\item the classical mutations make the freeze false and are rejected.
\end{itemize}

\section{P.5 entry dossier}\label{sec:p14-dossier}
One row per claim, each resolved mechanically by \texttt{dossier\_p14.py}: the function exists, the certificate key holds
the expected value, and the test exists. The five P.3/P.4 chains are carried over unchanged as frozen inputs; the P.5
entry chain is new.
{\scriptsize\setlength{\tabcolsep}{2.5pt}
\begin{longtable}{@{}p{0.9cm}p{4.1cm}p{3.5cm}p{5.0cm}p{3.0cm}@{}}
\toprule id & claim & file :: function & certificate key = value & test \\ \midrule \endhead
\multicolumn{5}{@{}l}{\textbf{P.3: projection -\textgreater{} dependency rows -\textgreater{} graphs [frozen P.3/P.4 input]}} \\
P3.1 & ch output projection fixed (c\^{}1 h\^{}1, ghost number 1, kappa\^{}2, one loop, transformed variables) & \texttt{ch\_\allowbreak{}source\_\allowbreak{}audit.\allowbreak{}py} :: \texttt{main} & \texttt{ch\_\allowbreak{}source\_\allowbreak{}audit/\allowbreak{}ch\_\allowbreak{}projection\_\allowbreak{}defined} = true\newline \texttt{ch\_\allowbreak{}source\_\allowbreak{}audit/\allowbreak{}ch\_\allowbreak{}pattern/\allowbreak{}ghost\_\allowbreak{}number} = 1 & \texttt{ch\_\allowbreak{}rows\_\allowbreak{}and\_\allowbreak{}quantum\_\allowbreak{}input\_\allowbreak{}slots} \\
P3.2 & vertex inventory generated from S\_Psi (not copied) & \texttt{ch\_\allowbreak{}source\_\allowbreak{}audit.\allowbreak{}py} :: \texttt{S\_\allowbreak{}psi\_\allowbreak{}inventory}, \texttt{vertex\_\allowbreak{}options} & \texttt{ch\_\allowbreak{}source\_\allowbreak{}audit/\allowbreak{}vertex\_\allowbreak{}inventory} = \{"kappa0": ["K.\allowbreak{}c",.\allowbreak{}.\allowbreak{}.\allowbreak{}\newline \texttt{ch\_\allowbreak{}source\_\allowbreak{}audit/\allowbreak{}checks/\allowbreak{}b\_\allowbreak{}has\_\allowbreak{}no\_\allowbreak{}interaction\_\allowbreak{}vertex} = true & \texttt{ch\_\allowbreak{}rows\_\allowbreak{}and\_\allowbreak{}quantum\_\allowbreak{}input\_\allowbreak{}slots} \\
P3.3 & dependency rows exhaustive: regenerated rows equal the published rows & \texttt{ch\_\allowbreak{}source\_\allowbreak{}audit.\allowbreak{}py} :: \texttt{audit}, \texttt{regenerate\_\allowbreak{}and\_\allowbreak{}compare} & \texttt{ch\_\allowbreak{}source\_\allowbreak{}audit/\allowbreak{}checks/\allowbreak{}regenerated\_\allowbreak{}rows\_\allowbreak{}match\_\allowbreak{}published} = true\newline \texttt{ch\_\allowbreak{}source\_\allowbreak{}audit/\allowbreak{}regeneration/\allowbreak{}complete} = true & \texttt{e2e\_\allowbreak{}dropped\_\allowbreak{}hstar\_\allowbreak{}c\_\allowbreak{}row} \\
P3.4 & 6 rows: 3 surviving, 3 derived zero (S2, cbar*, b* rows) & \texttt{ch\_\allowbreak{}source\_\allowbreak{}audit.\allowbreak{}py} :: \texttt{audit} & \texttt{ch\_\allowbreak{}source\_\allowbreak{}audit/\allowbreak{}rows\_\allowbreak{}total} = 6\newline \texttt{ch\_\allowbreak{}source\_\allowbreak{}audit/\allowbreak{}rows\_\allowbreak{}surviving} = 3\newline \texttt{ch\_\allowbreak{}source\_\allowbreak{}audit/\allowbreak{}rows\_\allowbreak{}derived\_\allowbreak{}zero} = 3 & \texttt{ch\_\allowbreak{}rows\_\allowbreak{}and\_\allowbreak{}quantum\_\allowbreak{}input\_\allowbreak{}slots} \\
P3.5 & enumeration lemma (justifies max\_len = 2) checked mechanically & \texttt{ch\_\allowbreak{}source\_\allowbreak{}audit.\allowbreak{}py} :: \texttt{enumeration\_\allowbreak{}lemma} & \texttt{ch\_\allowbreak{}source\_\allowbreak{}audit/\allowbreak{}enumeration\_\allowbreak{}lemma/\allowbreak{}holds} = true & \texttt{enumeration\_\allowbreak{}lemma} \\
P3.6 & graph topologies generated and named; no unnamed graph & \texttt{ch\_\allowbreak{}source\_\allowbreak{}audit.\allowbreak{}py} :: \texttt{enumerate\_\allowbreak{}graphs}, \texttt{name\_\allowbreak{}graph} & \texttt{graph\_\allowbreak{}inventory/\allowbreak{}names\_\allowbreak{}TBD} = []\newline \texttt{graph\_\allowbreak{}inventory/\allowbreak{}graphs} = \{"D03\_\allowbreak{}hh\_\allowbreak{}ghost\_\allowbreak{}loo.\allowbreak{}.\allowbreak{}.\allowbreak{}\newline \texttt{ch\_\allowbreak{}source\_\allowbreak{}audit/\allowbreak{}checks/\allowbreak{}names\_\allowbreak{}TBD\_\allowbreak{}empty} = true & \texttt{ch\_\allowbreak{}rows\_\allowbreak{}and\_\allowbreak{}quantum\_\allowbreak{}input\_\allowbreak{}slots} \\
P3.7 & P.3 flags & \texttt{release\_\allowbreak{}p14.\allowbreak{}py} :: \texttt{classical\_\allowbreak{}block} & \texttt{RELEASE\_\allowbreak{}STATUS/\allowbreak{}classical/\allowbreak{}flags/\allowbreak{}ch\_\allowbreak{}projection\_\allowbreak{}defined} = true\newline \texttt{RELEASE\_\allowbreak{}STATUS/\allowbreak{}classical/\allowbreak{}flags/\allowbreak{}ch\_\allowbreak{}dependency\_\allowbreak{}inventory\_\allowbreak{}complete} = true & \texttt{milestones\_\allowbreak{}consistent} \\
\midrule
\multicolumn{5}{@{}l}{\textbf{P.4: classical coefficient -\textgreater{} export -\textgreater{} verification [frozen P.3/P.4 input]}} \\
P4.A1 & C3 boundary identity sqrt(-g)R - L\_GG = -d\_mu B\^{}mu on generic jets & \texttt{boundary\_\allowbreak{}identity.\allowbreak{}py} :: \texttt{generic\_\allowbreak{}jet\_\allowbreak{}proof} & \texttt{boundary\_\allowbreak{}identity/\allowbreak{}C3\_\allowbreak{}proved} = true & \texttt{C3\_\allowbreak{}boundary\_\allowbreak{}identity\_\allowbreak{}generic\_\allowbreak{}and\_\allowbreak{}conformal} \\
P4.A2 & ell\_2, ell\_3, ell\_4 built and exported; export hash == quartic cert == master input & \texttt{quartic.\allowbreak{}py} :: \texttt{build\_\allowbreak{}and\_\allowbreak{}export} & \texttt{RELEASE\_\allowbreak{}STATUS/\allowbreak{}classical/\allowbreak{}hash\_\allowbreak{}consistency/\allowbreak{}consistent} = true\newline \texttt{quartic\_\allowbreak{}verified/\allowbreak{}ell4\_\allowbreak{}sha256} = "33ac9e49c4d1644ee.\allowbreak{}.\allowbreak{}.\allowbreak{} & \texttt{stale\_\allowbreak{}ell4\_\allowbreak{}hash\_\allowbreak{}blocks\_\allowbreak{}master} \\
P4.A3 & C1 inverse, C2 det/trace-log (10 symbolic components) & \texttt{quartic.\allowbreak{}py} :: \texttt{inverse\_\allowbreak{}control}, \texttt{det\_\allowbreak{}tracelog\_\allowbreak{}generic} & \texttt{quartic\_\allowbreak{}verified/\allowbreak{}acceptance\_\allowbreak{}controls\_\allowbreak{}deterministic/\allowbreak{}C1} = true\newline \texttt{quartic\_\allowbreak{}verified/\allowbreak{}acceptance\_\allowbreak{}controls\_\allowbreak{}deterministic/\allowbreak{}C2} = true & \texttt{quartic\_\allowbreak{}acceptance\_\allowbreak{}all\_\allowbreak{}deterministic} \\
P4.A4 & C4 export == independent evaluator (exact polynomial identity) & \texttt{quartic.\allowbreak{}py} :: \texttt{export\_\allowbreak{}identity} & \texttt{quartic\_\allowbreak{}verified/\allowbreak{}acceptance\_\allowbreak{}controls\_\allowbreak{}deterministic/\allowbreak{}C4} = true & \texttt{quartic\_\allowbreak{}acceptance\_\allowbreak{}all\_\allowbreak{}deterministic} \\
P4.A5 & C5/C6 ell\_3 = XII-C, ell\_2 = XI-A Fierz-Pauli modulo divergences (pinned XVI-N grav\_density) & \texttt{quartic.\allowbreak{}py} :: \texttt{regression\_\allowbreak{}deterministic}, \texttt{xvin\_\allowbreak{}symbolic} & \texttt{quartic\_\allowbreak{}verified/\allowbreak{}acceptance\_\allowbreak{}controls\_\allowbreak{}deterministic/\allowbreak{}C5} = true\newline \texttt{quartic\_\allowbreak{}verified/\allowbreak{}acceptance\_\allowbreak{}controls\_\allowbreak{}deterministic/\allowbreak{}C6} = true & \texttt{quartic\_\allowbreak{}acceptance\_\allowbreak{}all\_\allowbreak{}deterministic} \\
P4.A6 & C7 gravitational Noether identities on the exports, kappa\^{}0..kappa\^{}2 & \texttt{grav\_\allowbreak{}noether.\allowbreak{}py} :: \texttt{noether} & \texttt{quartic\_\allowbreak{}verified/\allowbreak{}controls/\allowbreak{}C7\_\allowbreak{}accepted\_\allowbreak{}orders} = \emph{all entries true} & \texttt{grav\_\allowbreak{}noether\_\allowbreak{}exact\_\allowbreak{}and\_\allowbreak{}mutation} \\
P4.A7 & quartic mutations rejected; quartic flag & \texttt{quartic.\allowbreak{}py} :: \texttt{main} & \texttt{quartic\_\allowbreak{}verified/\allowbreak{}mutations\_\allowbreak{}all\_\allowbreak{}rejected} = true\newline \texttt{RELEASE\_\allowbreak{}STATUS/\allowbreak{}classical/\allowbreak{}flags/\allowbreak{}quartic\_\allowbreak{}generator\_\allowbreak{}verified} = true & \texttt{e2e\_\allowbreak{}corrupted\_\allowbreak{}ell4\_\allowbreak{}coefficient} \\
\midrule
\multicolumn{5}{@{}l}{\textbf{P.4: transformations -\textgreater{} source terms [frozen P.3/P.4 input]}} \\
P4.B1 & ghost law derived from s\^{}2 = 0 (unique lambda = 1) & \texttt{brst\_\allowbreak{}from\_\allowbreak{}geometry.\allowbreak{}py} :: \texttt{solve\_\allowbreak{}ghost\_\allowbreak{}law} & \texttt{brst\_\allowbreak{}derivation/\allowbreak{}ghost\_\allowbreak{}law\_\allowbreak{}derivation/\allowbreak{}lam\_\allowbreak{}solutions} = ["1"] & \texttt{brst\_\allowbreak{}and\_\allowbreak{}source\_\allowbreak{}terms\_\allowbreak{}derived} \\
P4.B2 & s\_2 derived zero in every sector; s nilpotent at all orders & \texttt{brst\_\allowbreak{}from\_\allowbreak{}geometry.\allowbreak{}py} :: \texttt{per\_\allowbreak{}sector\_\allowbreak{}expansion}, \texttt{nilpotency} & \texttt{brst\_\allowbreak{}derivation/\allowbreak{}s2\_\allowbreak{}per\_\allowbreak{}sector} = \{"h": "derived\_\allowbreak{}zer.\allowbreak{}.\allowbreak{}.\allowbreak{}\newline \texttt{brst\_\allowbreak{}derivation/\allowbreak{}nilpotency\_\allowbreak{}exact\_\allowbreak{}all\_\allowbreak{}orders} = true & \texttt{brst\_\allowbreak{}and\_\allowbreak{}source\_\allowbreak{}terms\_\allowbreak{}derived} \\
P4.B3 & sigma\_A derived from the generator rule (unique); antifield-linear, no kappa\^{}2 X terms & \texttt{source\_\allowbreak{}action.\allowbreak{}py} :: \texttt{derive\_\allowbreak{}sigma}, \texttt{build\_\allowbreak{}X\_\allowbreak{}inv} & \texttt{source\_\allowbreak{}action/\allowbreak{}sigma\_\allowbreak{}unique} = true\newline \texttt{source\_\allowbreak{}action/\allowbreak{}X\_\allowbreak{}kappa2\_\allowbreak{}terms} = 0\newline \texttt{source\_\allowbreak{}action/\allowbreak{}antifield\_\allowbreak{}linear} = true & \texttt{brst\_\allowbreak{}and\_\allowbreak{}source\_\allowbreak{}terms\_\allowbreak{}derived}\newline \texttt{e2e\_\allowbreak{}flipped\_\allowbreak{}odd\_\allowbreak{}sign}\newline \texttt{e2e\_\allowbreak{}omitted\_\allowbreak{}cstar\_\allowbreak{}term} \\
\midrule
\multicolumn{5}{@{}l}{\textbf{P.4: gauge map -\textgreater{} Hessians -\textgreater{} source convention [frozen P.3/P.4 input]}} \\
P4.C1 & antifield shift orientation derived (eps = -1); S\_Psi = S + s Psi mod d & \texttt{gauge\_\allowbreak{}fixing\_\allowbreak{}map.\allowbreak{}py} :: \texttt{build}, \texttt{shift\_\allowbreak{}map} & \texttt{gauge\_\allowbreak{}fixing\_\allowbreak{}map/\allowbreak{}derived\_\allowbreak{}eps} = -1\newline \texttt{gauge\_\allowbreak{}fixing\_\allowbreak{}map/\allowbreak{}S\_\allowbreak{}Psi\_\allowbreak{}equals\_\allowbreak{}S\_\allowbreak{}plus\_\allowbreak{}sPsi\_\allowbreak{}mod\_\allowbreak{}d} = true & \texttt{e2e\_\allowbreak{}altered\_\allowbreak{}gauge\_\allowbreak{}shift} \\
P4.C2 & canonicity of the map on all 299 (anti)field pairs & \texttt{gauge\_\allowbreak{}fixing\_\allowbreak{}map.\allowbreak{}py} :: \texttt{canonicity} & \texttt{gauge\_\allowbreak{}fixing\_\allowbreak{}map/\allowbreak{}canonicity/\allowbreak{}pass} = true\newline \texttt{gauge\_\allowbreak{}fixing\_\allowbreak{}map/\allowbreak{}canonicity/\allowbreak{}pairs\_\allowbreak{}tested} = 299 & \texttt{gauge\_\allowbreak{}hessian\_\allowbreak{}generic\_\allowbreak{}and\_\allowbreak{}full\_\allowbreak{}canonicity} \\
P4.C3 & graviton Hessian at generic symbolic k: Schur complement k\^{}2 P, propagator (k\^{}2 P)\^{}-1 & \texttt{gauge\_\allowbreak{}fixing\_\allowbreak{}map.\allowbreak{}py} :: \texttt{hessian\_\allowbreak{}check} & \texttt{gauge\_\allowbreak{}fixing\_\allowbreak{}map/\allowbreak{}propagators/\allowbreak{}graviton/\allowbreak{}schur\_\allowbreak{}equals\_\allowbreak{}k2P} = true\newline \texttt{gauge\_\allowbreak{}fixing\_\allowbreak{}map/\allowbreak{}propagators/\allowbreak{}graviton/\allowbreak{}propagator\_\allowbreak{}hh\_\allowbreak{}equals\_\allowbreak{}inv\_\allowbreak{}k2P} = true & \texttt{gauge\_\allowbreak{}hessian\_\allowbreak{}generic\_\allowbreak{}and\_\allowbreak{}full\_\allowbreak{}canonicity} \\
P4.C4 & ghost (i cbar Box c) and scalar (-Box) kinetic terms; propagators match transformed Hessians & \texttt{gauge\_\allowbreak{}fixing\_\allowbreak{}map.\allowbreak{}py} :: \texttt{ghost\_\allowbreak{}kinetic\_\allowbreak{}check} & \texttt{gauge\_\allowbreak{}fixing\_\allowbreak{}map/\allowbreak{}propagators/\allowbreak{}ghost/\allowbreak{}match} = true\newline \texttt{gauge\_\allowbreak{}fixing\_\allowbreak{}map/\allowbreak{}propagators/\allowbreak{}scalar/\allowbreak{}match} = true\newline \texttt{gauge\_\allowbreak{}fixing\_\allowbreak{}map/\allowbreak{}propagators\_\allowbreak{}match\_\allowbreak{}transformed\_\allowbreak{}Hessians} = true & \texttt{gauge\_\allowbreak{}hessian\_\allowbreak{}generic\_\allowbreak{}and\_\allowbreak{}full\_\allowbreak{}canonicity} \\
P4.C5 & source convention Gamma\_Psi(Phi,Phi*') := Gamma(Phi, Phi*' - dPsi/dPhi): classical pullback, NOT a quantum theorem & \texttt{gauge\_\allowbreak{}fixing\_\allowbreak{}map.\allowbreak{}py} :: \texttt{build} & \texttt{gauge\_\allowbreak{}fixing\_\allowbreak{}map/\allowbreak{}source\_\allowbreak{}identity\_\allowbreak{}relation\_\allowbreak{}defined} = true\newline \texttt{gauge\_\allowbreak{}fixing\_\allowbreak{}map/\allowbreak{}source\_\allowbreak{}identity\_\allowbreak{}relation/\allowbreak{}status} = "DEFINED classical.\allowbreak{}.\allowbreak{}.\allowbreak{} & \texttt{gauge\_\allowbreak{}hessian\_\allowbreak{}generic\_\allowbreak{}and\_\allowbreak{}full\_\allowbreak{}canonicity} \\
P4.C6 & gauge-map flag & \texttt{release\_\allowbreak{}p14.\allowbreak{}py} :: \texttt{classical\_\allowbreak{}block} & \texttt{RELEASE\_\allowbreak{}STATUS/\allowbreak{}classical/\allowbreak{}flags/\allowbreak{}gauge\_\allowbreak{}fixing\_\allowbreak{}map\_\allowbreak{}constructed} = true & \texttt{milestones\_\allowbreak{}consistent} \\
\midrule
\multicolumn{5}{@{}l}{\textbf{P.4: master component -\textgreater{} residual record -\textgreater{} acceptance gate [frozen P.3/P.4 input]}} \\
P4.D1 & master equation at kappa\^{}0: every antifield-number component; grav af0 record; gate & \texttt{master\_\allowbreak{}kappa2.\allowbreak{}py} :: \texttt{run}, \texttt{XX\_\allowbreak{}by\_\allowbreak{}order}, \texttt{grav\_\allowbreak{}deterministic} & \texttt{master\_\allowbreak{}kappa2/\allowbreak{}table\_\allowbreak{}by\_\allowbreak{}kappa\_\allowbreak{}and\_\allowbreak{}antifield\_\allowbreak{}number/\allowbreak{}kappa0} = \emph{all entries true}\newline \texttt{master\_\allowbreak{}kappa2/\allowbreak{}details/\allowbreak{}grav\_\allowbreak{}af0\_\allowbreak{}deterministic/\allowbreak{}kappa0/\allowbreak{}Q\_\allowbreak{}terms} = [0, 0, 0, 0]\newline \texttt{master\_\allowbreak{}kappa2/\allowbreak{}details/\allowbreak{}grav\_\allowbreak{}af0\_\allowbreak{}deterministic/\allowbreak{}kappa0/\allowbreak{}Q\_\allowbreak{}identically\_\allowbreak{}zero} = true\newline \texttt{RELEASE\_\allowbreak{}STATUS/\allowbreak{}classical/\allowbreak{}residual\_\allowbreak{}records/\allowbreak{}master\_\allowbreak{}grav\_\allowbreak{}per\_\allowbreak{}order/\allowbreak{}kappa0/\allowbreak{}accepted} = true & \texttt{residual\_\allowbreak{}records\_\allowbreak{}regenerated\_\allowbreak{}and\_\allowbreak{}agree} \\
P4.D2 & master equation at kappa\^{}1: every antifield-number component; grav af0 record; gate & \texttt{master\_\allowbreak{}kappa2.\allowbreak{}py} :: \texttt{run}, \texttt{XX\_\allowbreak{}by\_\allowbreak{}order}, \texttt{grav\_\allowbreak{}deterministic} & \texttt{master\_\allowbreak{}kappa2/\allowbreak{}table\_\allowbreak{}by\_\allowbreak{}kappa\_\allowbreak{}and\_\allowbreak{}antifield\_\allowbreak{}number/\allowbreak{}kappa1} = \emph{all entries true}\newline \texttt{master\_\allowbreak{}kappa2/\allowbreak{}details/\allowbreak{}grav\_\allowbreak{}af0\_\allowbreak{}deterministic/\allowbreak{}kappa1/\allowbreak{}Q\_\allowbreak{}terms} = [0, 0, 0, 0]\newline \texttt{master\_\allowbreak{}kappa2/\allowbreak{}details/\allowbreak{}grav\_\allowbreak{}af0\_\allowbreak{}deterministic/\allowbreak{}kappa1/\allowbreak{}Q\_\allowbreak{}identically\_\allowbreak{}zero} = true\newline \texttt{RELEASE\_\allowbreak{}STATUS/\allowbreak{}classical/\allowbreak{}residual\_\allowbreak{}records/\allowbreak{}master\_\allowbreak{}grav\_\allowbreak{}per\_\allowbreak{}order/\allowbreak{}kappa1/\allowbreak{}accepted} = true & \texttt{residual\_\allowbreak{}records\_\allowbreak{}regenerated\_\allowbreak{}and\_\allowbreak{}agree} \\
P4.D3 & master equation at kappa\^{}2: every antifield-number component; grav af0 record; gate & \texttt{master\_\allowbreak{}kappa2.\allowbreak{}py} :: \texttt{run}, \texttt{XX\_\allowbreak{}by\_\allowbreak{}order}, \texttt{grav\_\allowbreak{}deterministic} & \texttt{master\_\allowbreak{}kappa2/\allowbreak{}table\_\allowbreak{}by\_\allowbreak{}kappa\_\allowbreak{}and\_\allowbreak{}antifield\_\allowbreak{}number/\allowbreak{}kappa2} = \emph{all entries true}\newline \texttt{master\_\allowbreak{}kappa2/\allowbreak{}details/\allowbreak{}grav\_\allowbreak{}af0\_\allowbreak{}deterministic/\allowbreak{}kappa2/\allowbreak{}Q\_\allowbreak{}terms} = [0, 0, 0, 0]\newline \texttt{master\_\allowbreak{}kappa2/\allowbreak{}details/\allowbreak{}grav\_\allowbreak{}af0\_\allowbreak{}deterministic/\allowbreak{}kappa2/\allowbreak{}Q\_\allowbreak{}identically\_\allowbreak{}zero} = true\newline \texttt{RELEASE\_\allowbreak{}STATUS/\allowbreak{}classical/\allowbreak{}residual\_\allowbreak{}records/\allowbreak{}master\_\allowbreak{}grav\_\allowbreak{}per\_\allowbreak{}order/\allowbreak{}kappa2/\allowbreak{}accepted} = true & \texttt{residual\_\allowbreak{}records\_\allowbreak{}regenerated\_\allowbreak{}and\_\allowbreak{}agree} \\
P4.D4 & no required master residual; no unconditional pass field & \texttt{master\_\allowbreak{}kappa2.\allowbreak{}py} :: \texttt{master\_\allowbreak{}residuals} & \texttt{master\_\allowbreak{}kappa2/\allowbreak{}required\_\allowbreak{}master\_\allowbreak{}residuals} = [] & \texttt{master\_\allowbreak{}no\_\allowbreak{}unconditional\_\allowbreak{}pass\_\allowbreak{}and\_\allowbreak{}deterministic} \\
P4.D5 & count/flag invariant on every record; quartic and master records agree & \texttt{certio.\allowbreak{}py} :: \texttt{invariant\_\allowbreak{}violations}, \texttt{cross\_\allowbreak{}check} & \texttt{RELEASE\_\allowbreak{}STATUS/\allowbreak{}classical/\allowbreak{}deterministic\_\allowbreak{}prerequisites/\allowbreak{}residual\_\allowbreak{}record\_\allowbreak{}invariants\_\allowbreak{}all\_\allowbreak{}certificates} = true\newline \texttt{RELEASE\_\allowbreak{}STATUS/\allowbreak{}classical/\allowbreak{}residual\_\allowbreak{}records/\allowbreak{}cross\_\allowbreak{}certificate/\allowbreak{}agree} = true & \texttt{certificate\_\allowbreak{}mutation\_\allowbreak{}rejects\_\allowbreak{}publication}\newline \texttt{certio\_\allowbreak{}rejects\_\allowbreak{}inconsistent\_\allowbreak{}record\_\allowbreak{}before\_\allowbreak{}serialization} \\
P4.D6 & acceptance gate: release\_claim\_accepted and classical readiness (P.4) & \texttt{release\_\allowbreak{}p14.\allowbreak{}py} :: \texttt{classical\_\allowbreak{}block} & \texttt{RELEASE\_\allowbreak{}STATUS/\allowbreak{}classical/\allowbreak{}release\_\allowbreak{}claim\_\allowbreak{}accepted} = true\newline \texttt{RELEASE\_\allowbreak{}STATUS/\allowbreak{}classical/\allowbreak{}classical\_\allowbreak{}readiness\_\allowbreak{}for\_\allowbreak{}ch\_\allowbreak{}identity} = true\newline \texttt{RELEASE\_\allowbreak{}STATUS/\allowbreak{}classical/\allowbreak{}flags/\allowbreak{}required\_\allowbreak{}classical\_\allowbreak{}BV\_\allowbreak{}components\_\allowbreak{}verified} = true & \texttt{release\_\allowbreak{}accepted\_\allowbreak{}and\_\allowbreak{}hard\_\allowbreak{}stops} \\
P4.D7 & publication gate re-invoked before every pdf/zip/manifest stage, also on resumed builds & \texttt{build.\allowbreak{}py} :: \texttt{publication\_\allowbreak{}gate}, \texttt{pdf\_\allowbreak{}readback} & (build-time gate) & \texttt{resumed\_\allowbreak{}stage\_\allowbreak{}build\_\allowbreak{}blocked} \\
\midrule
\multicolumn{5}{@{}l}{\textbf{P.5\_entry: quantum object -\textgreater{} prescription/D11 matching -\textgreater{} counterterm domain -\textgreater{} amplitude rules -\textgreater{} gate}} \\
P5E.1 & quantum object in transformed variables: S\_Psi parts, eps = -1, b retained (elimination makes S quadratic in cbar*'), Gamma\_1 independent of cbar*', b*'; pullback != quantum theorem & \texttt{quantum\_\allowbreak{}object.\allowbreak{}py} :: \texttt{main}, \texttt{patterns}, \texttt{b\_\allowbreak{}elimination} & \texttt{\allowbreak{}quantum\_\allowbreak{}object/\allowbreak{}quantum\_\allowbreak{}action\_\allowbreak{}defined} = true\newline \texttt{\allowbreak{}quantum\_\allowbreak{}object/\allowbreak{}Q1\_\allowbreak{}transformed\_\allowbreak{}variables/\allowbreak{}eps\_\allowbreak{}derived} = -1\newline \texttt{\allowbreak{}quantum\_\allowbreak{}object/\allowbreak{}Q2\_\allowbreak{}fields\_\allowbreak{}antifields\_\allowbreak{}b/\allowbreak{}b\_\allowbreak{}decision} = "RETAINED and integrated"\newline \texttt{\allowbreak{}quantum\_\allowbreak{}object/\allowbreak{}Q2\_\allowbreak{}fields\_\allowbreak{}antifields\_\allowbreak{}b/\allowbreak{}Gamma1\_\allowbreak{}independent\_\allowbreak{}of} = ["bs", "cbs"]\newline \texttt{\allowbreak{}quantum\_\allowbreak{}object/\allowbreak{}Q5\_\allowbreak{}kappa\_\allowbreak{}loop\_\allowbreak{}bookkeeping/\allowbreak{}ok} = true\newline \texttt{\allowbreak{}quantum\_\allowbreak{}object/\allowbreak{}Q6\_\allowbreak{}pullback\_\allowbreak{}vs\_\allowbreak{}theorem/\allowbreak{}quantum\_\allowbreak{}gauge\_\allowbreak{}change\_\allowbreak{}theorem\_\allowbreak{}asserted} = false & \texttt{quantum\_\allowbreak{}object} \\
P5E.2 & ONE prescription: operational DRED (rules R1-R6), MS-bar, Feynman i0; scaleless rules defined but not applied; quantum\_prescription\_defined requires the verified DRED rules & \texttt{prescription\_\allowbreak{}d11.\allowbreak{}py} :: \texttt{main} & \texttt{\allowbreak{}prescription\_\allowbreak{}D11\_\allowbreak{}matching/\allowbreak{}quantum\_\allowbreak{}prescription\_\allowbreak{}defined} = true\newline \texttt{\allowbreak{}prescription\_\allowbreak{}D11\_\allowbreak{}matching/\allowbreak{}prescription/\allowbreak{}id} = "\allowbreak{}DRED\_\allowbreak{}MSbar\_\allowbreak{}FeynmanI0"\newline \texttt{\allowbreak{}prescription\_\allowbreak{}D11\_\allowbreak{}matching/\allowbreak{}prescription\_\allowbreak{}checks} = \emph{all entries true}\newline \texttt{\allowbreak{}prescription\_\allowbreak{}D11\_\allowbreak{}matching/\allowbreak{}prescription\_\allowbreak{}checks/\allowbreak{}operational\_\allowbreak{}DRED\_\allowbreak{}rules\_\allowbreak{}verified} = true\newline \texttt{\allowbreak{}prescription\_\allowbreak{}D11\_\allowbreak{}matching/\allowbreak{}D12\_\allowbreak{}status} = "Missing" & \texttt{prescription\_\allowbreak{}and\_\allowbreak{}D11\_\allowbreak{}matching} \\
P5E.2b & operational DRED: g4 = ghat + gtilde (tr d, 4-d), l and p in QdS, l\_r = l.e\_(r) lifting, ghat-projected Gram/PV reduction, contractions before Laurent with d symbolic; evanescent parts retained; premature d=4, tr ghat=4 and expand-before-contract are rejected & \texttt{dred.\allowbreak{}py} :: \texttt{reduce\_\allowbreak{}tensor}, \texttt{eta\_\allowbreak{}trace\_\allowbreak{}Gt}, \texttt{toy\_\allowbreak{}four\_\allowbreak{}d\_\allowbreak{}part}, \texttt{toy\_\allowbreak{}evanescent\_\allowbreak{}trace}, \texttt{known\_\allowbreak{}reductions}, \texttt{lift\_\allowbreak{}exported}, \texttt{main} & \texttt{\allowbreak{}dred\_\allowbreak{}rules/\allowbreak{}dred\_\allowbreak{}rules\_\allowbreak{}verified} = true\newline \texttt{\allowbreak{}dred\_\allowbreak{}rules/\allowbreak{}mutations\_\allowbreak{}all\_\allowbreak{}rejected} = true\newline \texttt{\allowbreak{}dred\_\allowbreak{}rules/\allowbreak{}mutations/\allowbreak{}premature\_\allowbreak{}d4/\allowbreak{}rejected} = true\newline \texttt{\allowbreak{}dred\_\allowbreak{}rules/\allowbreak{}mutations/\allowbreak{}trace\_\allowbreak{}ghat\_\allowbreak{}eq\_\allowbreak{}4/\allowbreak{}rejected} = true\newline \texttt{\allowbreak{}dred\_\allowbreak{}rules/\allowbreak{}mutations/\allowbreak{}contract\_\allowbreak{}after\_\allowbreak{}laurent/\allowbreak{}rejected} = true\newline \texttt{\allowbreak{}dred\_\allowbreak{}rules/\allowbreak{}toy\_\allowbreak{}checks/\allowbreak{}T1\_\allowbreak{}four\_\allowbreak{}d\_\allowbreak{}part/\allowbreak{}finite} = "a/\allowbreak{}8 + b/\allowbreak{}4"\newline \texttt{\allowbreak{}dred\_\allowbreak{}rules/\allowbreak{}toy\_\allowbreak{}checks/\allowbreak{}T2\_\allowbreak{}evanescent\_\allowbreak{}trace/\allowbreak{}finite} = "-a/\allowbreak{}2"\newline \texttt{\allowbreak{}dred\_\allowbreak{}rules/\allowbreak{}known\_\allowbreak{}reductions/\allowbreak{}pass} = true\newline \texttt{\allowbreak{}dred\_\allowbreak{}rules/\allowbreak{}evanescent\_\allowbreak{}parts\_\allowbreak{}present\_\allowbreak{}in\_\allowbreak{}exports} = true\newline \texttt{\allowbreak{}dred\_\allowbreak{}rules/\allowbreak{}loop\_\allowbreak{}evaluations\_\allowbreak{}performed} = false & \texttt{dred\_\allowbreak{}rules\_\allowbreak{}certificate}\newline \texttt{dred\_\allowbreak{}premature\_\allowbreak{}d4\_\allowbreak{}changes\_\allowbreak{}finite\_\allowbreak{}part\_\allowbreak{}and\_\allowbreak{}is\_\allowbreak{}caught}\newline \texttt{dred\_\allowbreak{}evanescent\_\allowbreak{}trace\_\allowbreak{}is\_\allowbreak{}d\_\allowbreak{}minus\_\allowbreak{}4\_\allowbreak{}times\_\allowbreak{}pole}\newline \texttt{dred\_\allowbreak{}trace\_\allowbreak{}ghat\_\allowbreak{}eq\_\allowbreak{}4\_\allowbreak{}rejected}\newline \texttt{dred\_\allowbreak{}known\_\allowbreak{}reductions\_\allowbreak{}and\_\allowbreak{}contraction\_\allowbreak{}engine} \\
P5E.3 & D11 matching DEFINED (M1 exact integrand identity vs XVI-O V\_tensor, 100 pairs; M3 CONDITIONAL on M2: sd(u\^{}R) 4, kernel sd \textless{}= 8, contact derivative order \textless{}= 4, contact space orders 0/2/4 = 2/4/5, degree 4 imposed by normalization condition N4; no c delta; r = -i pi\^{}2 XVI-O side only); NOT verified & \texttt{prescription\_\allowbreak{}d11.\allowbreak{}py} :: \texttt{M1\_\allowbreak{}identity}, \texttt{M3\_\allowbreak{}local\_\allowbreak{}space}, \texttt{M3\_\allowbreak{}contact\_\allowbreak{}analysis} & \texttt{\allowbreak{}prescription\_\allowbreak{}D11\_\allowbreak{}matching/\allowbreak{}D11\_\allowbreak{}matching\_\allowbreak{}defined} = true\newline \texttt{\allowbreak{}prescription\_\allowbreak{}D11\_\allowbreak{}matching/\allowbreak{}D11\_\allowbreak{}matching\_\allowbreak{}verified} = false\newline \texttt{\allowbreak{}prescription\_\allowbreak{}D11\_\allowbreak{}matching/\allowbreak{}matching/\allowbreak{}M1\_\allowbreak{}integrand\_\allowbreak{}identity/\allowbreak{}mismatches} = []\newline \texttt{\allowbreak{}prescription\_\allowbreak{}D11\_\allowbreak{}matching/\allowbreak{}matching/\allowbreak{}M1\_\allowbreak{}integrand\_\allowbreak{}identity/\allowbreak{}component\_\allowbreak{}pairs\_\allowbreak{}checked} = 100\newline \texttt{\allowbreak{}prescription\_\allowbreak{}D11\_\allowbreak{}matching/\allowbreak{}matching/\allowbreak{}M1\_\allowbreak{}integrand\_\allowbreak{}identity/\allowbreak{}bound\_\allowbreak{}to\_\allowbreak{}export} = true\newline \texttt{\allowbreak{}prescription\_\allowbreak{}D11\_\allowbreak{}matching/\allowbreak{}matching/\allowbreak{}M3\_\allowbreak{}local\_\allowbreak{}difference/\allowbreak{}labels/\allowbreak{}scalar\_\allowbreak{}extension\_\allowbreak{}scaling\_\allowbreak{}degree} = 4\newline \texttt{\allowbreak{}prescription\_\allowbreak{}D11\_\allowbreak{}matching/\allowbreak{}matching/\allowbreak{}M3\_\allowbreak{}local\_\allowbreak{}difference/\allowbreak{}labels/\allowbreak{}kernel\_\allowbreak{}scaling\_\allowbreak{}degree\_\allowbreak{}bound} = 8\newline \texttt{\allowbreak{}prescription\_\allowbreak{}D11\_\allowbreak{}matching/\allowbreak{}matching/\allowbreak{}M3\_\allowbreak{}local\_\allowbreak{}difference/\allowbreak{}labels/\allowbreak{}local\_\allowbreak{}contact\_\allowbreak{}derivative\_\allowbreak{}order\_\allowbreak{}bound} = 4\newline \texttt{\allowbreak{}prescription\_\allowbreak{}D11\_\allowbreak{}matching/\allowbreak{}matching/\allowbreak{}M3\_\allowbreak{}local\_\allowbreak{}difference/\allowbreak{}lower\_\allowbreak{}degree\_\allowbreak{}contacts\_\allowbreak{}excluded\_\allowbreak{}by\_\allowbreak{}N4/\allowbreak{}order\_\allowbreak{}2\_\allowbreak{}mu2} = 4\newline \texttt{\allowbreak{}prescription\_\allowbreak{}D11\_\allowbreak{}matching/\allowbreak{}matching/\allowbreak{}M3\_\allowbreak{}local\_\allowbreak{}difference/\allowbreak{}dimension} = 5\newline \texttt{\allowbreak{}prescription\_\allowbreak{}D11\_\allowbreak{}matching/\allowbreak{}matching/\allowbreak{}M3\_\allowbreak{}local\_\allowbreak{}difference/\allowbreak{}transverse\_\allowbreak{}subspace\_\allowbreak{}dimension} = 2\newline \texttt{\allowbreak{}prescription\_\allowbreak{}D11\_\allowbreak{}matching/\allowbreak{}matching/\allowbreak{}M4\_\allowbreak{}finite\_\allowbreak{}adjustments/\allowbreak{}u\_\allowbreak{}R\_\allowbreak{}c\_\allowbreak{}delta\_\allowbreak{}added\_\allowbreak{}by\_\allowbreak{}P} = false\newline \texttt{\allowbreak{}prescription\_\allowbreak{}D11\_\allowbreak{}matching/\allowbreak{}matching/\allowbreak{}M5\_\allowbreak{}residue/\allowbreak{}r} = "-i * pi**2" & \texttt{prescription\_\allowbreak{}and\_\allowbreak{}D11\_\allowbreak{}matching} \\
P5E.4 & counterterm domain: kappa\^{}1 empty, kappa\^{}2 basis 11 (hh 5, cbar c 2, K'c 3, phi phi 1) generated mechanically; eight candidate directions modulo the computed exact equivalences, before imposing the identity constraints (no cohomological classification: D cap ker sigma\_0 not computed); constraint list, not an admission & \texttt{counterterm\_\allowbreak{}domain.\allowbreak{}py} :: \texttt{structures}, \texttt{admissible\_\allowbreak{}pairs}, \texttt{rank}, \texttt{main} & \texttt{\allowbreak{}counterterm\_\allowbreak{}domain/\allowbreak{}counterterm\_\allowbreak{}domain\_\allowbreak{}defined} = true\newline \texttt{\allowbreak{}counterterm\_\allowbreak{}domain/\allowbreak{}kappa1/\allowbreak{}basis\_\allowbreak{}size} = 0\newline \texttt{\allowbreak{}counterterm\_\allowbreak{}domain/\allowbreak{}kappa2/\allowbreak{}basis\_\allowbreak{}size} = 11\newline \texttt{\allowbreak{}counterterm\_\allowbreak{}domain/\allowbreak{}equivalences/\allowbreak{}quotient\_\allowbreak{}dimension} = 8\newline \texttt{\allowbreak{}counterterm\_\allowbreak{}domain/\allowbreak{}equivalences/\allowbreak{}D\_\allowbreak{}cap\_\allowbreak{}ker\_\allowbreak{}sigma0\_\allowbreak{}computed} = false\newline \texttt{\allowbreak{}counterterm\_\allowbreak{}domain/\allowbreak{}equivalences/\allowbreak{}S0\_\allowbreak{}density\_\allowbreak{}check/\allowbreak{}ok} = true\newline \texttt{\allowbreak{}counterterm\_\allowbreak{}domain/\allowbreak{}XIV\_\allowbreak{}J\_\allowbreak{}opened} = false\newline \texttt{\allowbreak{}counterterm\_\allowbreak{}domain/\allowbreak{}XV\_\allowbreak{}B\_\allowbreak{}opened} = false & \texttt{counterterm\_\allowbreak{}domain} \\
P5E.5 & amplitude rules: vertices (incl. ghost-graviton derivative terms, K' vertex), propagators, ordering, multiplicities, symmetry factors, signs; symbolic integrands; W1-W6 controls (W5 regenerates only D14a-c, D11, D13; W6 regenerates D01, D02, D03, D12) & \texttt{amplitude\_\allowbreak{}rules.\allowbreak{}py} :: \texttt{build\_\allowbreak{}integrands}, \texttt{verify}, \texttt{W2}, \texttt{W3}, \texttt{W4}, \texttt{W5} & \texttt{\allowbreak{}amplitude\_\allowbreak{}rules/\allowbreak{}amplitude\_\allowbreak{}rules\_\allowbreak{}verified} = true\newline \texttt{\allowbreak{}amplitude\_\allowbreak{}rules/\allowbreak{}controls/\allowbreak{}W1\_\allowbreak{}wick\_\allowbreak{}berezin\_\allowbreak{}shadow/\allowbreak{}pass} = true\newline \texttt{\allowbreak{}amplitude\_\allowbreak{}rules/\allowbreak{}controls/\allowbreak{}W4\_\allowbreak{}ward\_\allowbreak{}bose/\allowbreak{}pass} = true\newline \texttt{\allowbreak{}amplitude\_\allowbreak{}rules/\allowbreak{}integrands\_\allowbreak{}only} = true\newline \texttt{\allowbreak{}amplitude\_\allowbreak{}rules/\allowbreak{}loop\_\allowbreak{}evaluations\_\allowbreak{}performed} = false\newline \texttt{\allowbreak{}amplitude\_\allowbreak{}rules/\allowbreak{}graphs/\allowbreak{}D12/\allowbreak{}D12\_\allowbreak{}status} = "Missing" & \texttt{amplitude\_\allowbreak{}rules\_\allowbreak{}and\_\allowbreak{}mutations} \\
P5E.5b & independent Hessian/Tr-log regeneration of D01, D02, D03, D12 from the field-dependent Hessian of S\_Psi (own (h,b) inverse; not the Feynman-rule assembler) == exported integrands & \texttt{hessian\_\allowbreak{}route.\allowbreak{}py} :: \texttt{ell\_\allowbreak{}table}, \texttt{graviton\_\allowbreak{}inverse}, \texttt{ghost\_\allowbreak{}inverse}, \texttt{reference}, \texttt{W6}, \texttt{compare\_\allowbreak{}exports} & \texttt{\allowbreak{}amplitude\_\allowbreak{}rules/\allowbreak{}controls/\allowbreak{}W6\_\allowbreak{}hessian\_\allowbreak{}route/\allowbreak{}pass} = true\newline \texttt{\allowbreak{}amplitude\_\allowbreak{}rules/\allowbreak{}W6\_\allowbreak{}exports/\allowbreak{}all\_\allowbreak{}agree} = true\newline \texttt{\allowbreak{}amplitude\_\allowbreak{}rules/\allowbreak{}W6\_\allowbreak{}exports/\allowbreak{}D01/\allowbreak{}n\_\allowbreak{}mismatches} = 0\newline \texttt{\allowbreak{}amplitude\_\allowbreak{}rules/\allowbreak{}W6\_\allowbreak{}exports/\allowbreak{}D03/\allowbreak{}n\_\allowbreak{}mismatches} = 0 & \texttt{hessian\_\allowbreak{}route\_\allowbreak{}regenerates\_\allowbreak{}D01\_\allowbreak{}D02\_\allowbreak{}D03\_\allowbreak{}D12} \\
P5E.6 & mutations rejected: tadpole factor, ghost-loop sign, off-diagonal multiplicity, source/ghost order, ghost transport term & \texttt{amplitude\_\allowbreak{}rules.\allowbreak{}py} :: \texttt{apply\_\allowbreak{}mutations}, \texttt{main} & \texttt{\allowbreak{}amplitude\_\allowbreak{}rules/\allowbreak{}mutations\_\allowbreak{}all\_\allowbreak{}rejected} = true\newline \texttt{\allowbreak{}amplitude\_\allowbreak{}rules/\allowbreak{}mutations/\allowbreak{}omit\_\allowbreak{}ghost\_\allowbreak{}transport\_\allowbreak{}term/\allowbreak{}rejected} = true\newline \texttt{\allowbreak{}amplitude\_\allowbreak{}rules/\allowbreak{}mutations/\allowbreak{}reversed\_\allowbreak{}source\_\allowbreak{}ghost\_\allowbreak{}order/\allowbreak{}rejected} = true & \texttt{amplitude\_\allowbreak{}rules\_\allowbreak{}and\_\allowbreak{}mutations} \\
P5E.7 & Berezin/Wick shadow fixes symmetry factors and Grassmann signs independently & \texttt{shadow.\allowbreak{}py} :: \texttt{expect}, \texttt{gamma1}, \texttt{run\_\allowbreak{}shadows} & \texttt{\allowbreak{}amplitude\_\allowbreak{}rules/\allowbreak{}controls/\allowbreak{}W1\_\allowbreak{}wick\_\allowbreak{}berezin\_\allowbreak{}shadow/\allowbreak{}pass} = true & \texttt{amplitude\_\allowbreak{}rules\_\allowbreak{}and\_\allowbreak{}mutations} \\
P5E.8 & P.3/P.4 freeze bound to the p13 tree (independent rebuild passed; same-box caveat); frozen code and values unchanged & \texttt{freeze\_\allowbreak{}p13.\allowbreak{}py} :: \texttt{freeze}, \texttt{F1\_\allowbreak{}independent}, \texttt{F2\_\allowbreak{}frozen\_\allowbreak{}code}, \texttt{F3\_\allowbreak{}frozen\_\allowbreak{}values} & \texttt{RELEASE\_\allowbreak{}STATUS/\allowbreak{}P3\_\allowbreak{}P4\_\allowbreak{}frozen\_\allowbreak{}as\_\allowbreak{}P5\_\allowbreak{}inputs} = true\newline \texttt{RELEASE\_\allowbreak{}STATUS/\allowbreak{}p13\_\allowbreak{}independent\_\allowbreak{}rebuild/\allowbreak{}bound\_\allowbreak{}to\_\allowbreak{}source\_\allowbreak{}tree\_\allowbreak{}sha256} = "2b1baf335468944aa.\allowbreak{}.\allowbreak{}.\allowbreak{}\newline \texttt{RELEASE\_\allowbreak{}STATUS/\allowbreak{}p13\_\allowbreak{}independent\_\allowbreak{}rebuild/\allowbreak{}caveat} = "the p13 independe.\allowbreak{}.\allowbreak{}.\allowbreak{}\newline \texttt{RELEASE\_\allowbreak{}STATUS/\allowbreak{}freeze/\allowbreak{}F2/\allowbreak{}changed\_\allowbreak{}or\_\allowbreak{}missing} = []\newline \texttt{RELEASE\_\allowbreak{}STATUS/\allowbreak{}freeze/\allowbreak{}F3/\allowbreak{}differences} = [] & \texttt{freeze\_\allowbreak{}binding} \\
P5E.9 & P.5 entry gate: earned iff all five contract flags hold; loop/restoration flags NOT earned & \texttt{release\_\allowbreak{}p14.\allowbreak{}py} :: \texttt{\_\allowbreak{}main}, \texttt{p5\_\allowbreak{}entry\_\allowbreak{}flags}, \texttt{milestones\_\allowbreak{}p14} & \texttt{RELEASE\_\allowbreak{}STATUS/\allowbreak{}P5\_\allowbreak{}entry\_\allowbreak{}flags} = \emph{all entries true}\newline \texttt{RELEASE\_\allowbreak{}STATUS/\allowbreak{}P5\_\allowbreak{}entry\_\allowbreak{}earned} = true\newline \texttt{RELEASE\_\allowbreak{}STATUS/\allowbreak{}hard\_\allowbreak{}stops\_\allowbreak{}hold} = true\newline \texttt{RELEASE\_\allowbreak{}STATUS/\allowbreak{}hard\_\allowbreak{}stops/\allowbreak{}quantum\_\allowbreak{}inputs\_\allowbreak{}evaluated} = "not\_\allowbreak{}earned"\newline \texttt{RELEASE\_\allowbreak{}STATUS/\allowbreak{}hard\_\allowbreak{}stops/\allowbreak{}projected\_\allowbreak{}identity\_\allowbreak{}verified} = "not\_\allowbreak{}earned"\newline \texttt{RELEASE\_\allowbreak{}STATUS/\allowbreak{}hard\_\allowbreak{}stops/\allowbreak{}bulk\_\allowbreak{}symmetry\_\allowbreak{}restored} = false\newline \texttt{RELEASE\_\allowbreak{}STATUS/\allowbreak{}D11\_\allowbreak{}matching\_\allowbreak{}verified} = false\newline \texttt{RELEASE\_\allowbreak{}STATUS/\allowbreak{}D12\_\allowbreak{}status} = "Missing"\newline \texttt{RELEASE\_\allowbreak{}STATUS/\allowbreak{}hard\_\allowbreak{}stops/\allowbreak{}O1\_\allowbreak{}closure} = "NOT\_\allowbreak{}CLOSED" & \texttt{p5\_\allowbreak{}entry\_\allowbreak{}gate\_\allowbreak{}and\_\allowbreak{}hard\_\allowbreak{}stops} \\
P5E.10 & publication gate (p14) re-invoked before pdf/zip/manifest stages & \texttt{build.\allowbreak{}py} :: \texttt{publication\_\allowbreak{}gate}, \texttt{pdf\_\allowbreak{}readback} & (build-time gate) & \texttt{resumed\_\allowbreak{}stage\_\allowbreak{}build\_\allowbreak{}blocked} \\
\midrule
\multicolumn{5}{@{}l}{\textbf{Quantum input slots (named dependencies, NOT evaluated coefficients)}} \\
slot & hh at kappa\^{}2 ($\Gamma_1$ pattern \texttt{h\^{}2}) & \texttt{ch\_source\_audit.py} :: \texttt{audit} & \texttt{ch\_source\_audit/\allowbreak{}surviving\_\allowbreak{}rows/\allowbreak{}0} = D01, D02, D03, D11, D12 & \texttt{ch\_\allowbreak{}rows\_\allowbreak{}and\_\allowbreak{}quantum\_\allowbreak{}input\_\allowbreak{}slots} \\
slot & K'c at kappa\^{}2 ($\Gamma_1$ pattern \texttt{K.c}) & \texttt{ch\_source\_audit.py} :: \texttt{audit} & \texttt{ch\_source\_audit/\allowbreak{}surviving\_\allowbreak{}rows/\allowbreak{}1} = D13 & \texttt{ch\_\allowbreak{}rows\_\allowbreak{}and\_\allowbreak{}quantum\_\allowbreak{}input\_\allowbreak{}slots} \\
slot & h one-point at kappa ($\Gamma_1$ pattern \texttt{h}) & \texttt{ch\_source\_audit.py} :: \texttt{audit} & \texttt{ch\_source\_audit/\allowbreak{}surviving\_\allowbreak{}rows/\allowbreak{}2} = D14a, D14b, D14c & \texttt{ch\_\allowbreak{}rows\_\allowbreak{}and\_\allowbreak{}quantum\_\allowbreak{}input\_\allowbreak{}slots} \\
\midrule \multicolumn{5}{@{}p{16.5cm}@{}}{\textbf{Reproducibility flags (package evidence, read from RELEASE\_STATUS at build time; p14 flags bound to the p14 tree):} \texttt{author\_rebuild\_passed} = false; \texttt{independent\_rebuild\_passed} = false; \texttt{P3\_P4\_frozen\_as\_P5\_inputs} = true; p13 \texttt{independent\_rebuild\_passed} = true (bound to the p13 source tree \texttt{2b1baf335468944a...}; caveat: the p13 independent rebuild ran on the SAME box (shared kernel, TeX Live, poppler) in a fresh venv from the source-only zip; it was NOT run on a second machine).} \\
\bottomrule
\end{longtable}}


\section{Reproducibility contract (p14r)}\label{sec:p14-repro}
\begin{itemize}
\item \texttt{author\_rebuild\_passed} is bound to the \emph{p14r} source tree
(\texttt{reproducibility/author\_rebuild\_report\_p14r.json}).
\item The p14r \texttt{independent\_rebuild\_passed} is \texttt{pending}.
\item The p13 independent result is recorded separately and bound to the p13 tree; this is the same-box caveat above.
\item The source-only archive is \texttt{XVI-P-p14r-source-only.zip}, with reference values in
\texttt{reproducibility/reference\_values\_p14r.json} and entrypoint \texttt{python3 tools/rebuild\_from\_source.py}.
\item The p13 source-only zip, its reference values and the independent report are archived unchanged in
\texttt{archival/p13/}. The p14 release (zip, reference values, author report) is archived in \texttt{archival/p14/}.
\end{itemize}

\begin{nonclaim}
No loop is evaluated, no one-loop coefficient is fixed, and no symmetry is restored. P.5 entry fixes a contract: an
action, a prescription, a matching \emph{definition}, a counterterm \emph{domain} and integrands. P.5 calculation and
P.5 identity are open. The prepared $S_D$ keeps $(S_D,S_D)\neq0$ (XVI-N prop:broken-square), and no completion
contradicting XVI-N thm:massless-obstruction is claimed. En01, Em01, Em02-B/D, En02, VII-E and VIII-B are not premises.
\end{nonclaim}
